%==============================================================
%  lecons/ensembles/construction_H.tex — Construction de H
%==============================================================

\lecontitle{Construction de \texorpdfstring{$\mathbb{H}$}{H}}{Théorie des ensembles — Quaternions de Hamilton}

%=============================================================
%  § 1 — MOTIVATION ET DÉFINITION
%=============================================================
\secnum{navyblue}{1}{Définition des quaternions}

\rubriq{Contexte historique}
Hamilton cherchait (1843) à étendre $\mathbb{C}$ pour représenter les rotations
de l'espace $\mathbb{R}^3$. Après des années à tenter une structure en
dimension $3$, il découvrit que la dimension $4$ est nécessaire, en renonçant
à la commutativité de la multiplication.

\begin{tcolorbox}[thmbox]
  \textbf{Définition.}\enspace%
  \index{quaternions!définition}\index{Hamilton (quaternions de)}
  L'algèbre des \emph{quaternions} est
  \[
    \mathbb{H} = \{a + bi + cj + dk \mid a,b,c,d \in \mathbb{R}\}
  \]
  munie de l'addition composante par composante et de la multiplication
  déterminée par les relations de Hamilton :
  \[
    i^2 = j^2 = k^2 = ijk = -1.
  \]
  Ces relations imposent :
  $ij = k,\ ji = -k,\ jk = i,\ kj = -i,\ ki = j,\ ik = -j$.
\end{tcolorbox}

\decsep{}

%=============================================================
%  § 2 — STRUCTURE ALGÉBRIQUE
%=============================================================
\secnum{navyblue}{2}{$\mathbb{H}$ est un corps gauche}

\begin{tcolorbox}[thmbox]
  \textbf{Propriétés fondamentales.}\enspace%
  \index{quaternions!norme}\index{quaternions!conjugué}\index{corps gauche}
  Pour $q = a + bi + cj + dk \in \mathbb{H}$, on définit :
  \[
    \bar{q} = a - bi - cj - dk \quad \text{(conjugué)},
    \qquad
    |q|^2 = q\bar{q} = a^2+b^2+c^2+d^2 \quad \text{(norme au carré)}.
  \]
  \begin{enumerate}
    \item $\overline{pq} = \bar{q}\bar{p}$ pour tous $p,q \in \mathbb{H}$.
    \item $|pq| = |p||q|$ (multiplicativité de la norme).
    \item Pour $q \neq 0$ : $q^{-1} = \bar{q}/|q|^2$.
    \item $(\mathbb{H}, +, \times)$ est un \emph{corps gauche} (division ring) :
          tous les axiomes d'un corps sauf la commutativité de $\times$.
  \end{enumerate}
\end{tcolorbox}

\begin{mdframed}[style=proofbar]
  \textit{Inverse.} $q \cdot \bar{q} = |q|^2 \in \mathbb{R}_+^*$, donc
  $q \cdot (\bar{q}/|q|^2) = 1$. De même $(\bar{q}/|q|^2) \cdot q = 1$,
  ce qui donne bien un inverse à gauche et à droite.

  \textit{Associativité.} Vérification directe (fastidieuse mais mécanique) :
  on vérifie les $64$ triplets de base ${\{1,i,j,k\}}^3$.

  \textit{Non-commutativité.} $ij = k \neq -k = ji$.
  \hfill$\blacksquare$
\end{mdframed}

\rubriq{Construction matricielle}

\begin{mdframed}[style=proofbar]
  $\mathbb{H}$ est isomorphe (comme algèbre réelle) au sous-anneau de
  $\mathcal{M}_2(\mathbb{C})$ :
  \[
    a + bi + cj + dk \;\longmapsto\;
    \begin{pmatrix} a+bi & c+di \\ -(c-di) & a-bi \end{pmatrix}
    =
    \begin{pmatrix} a+bi & c+di \\ -\overline{c+di} & \overline{a+bi} \end{pmatrix}.
  \]
  Pour établir que l'application est un morphisme d'algèbres, il suffit de
  vérifier que les images de $1,i,j,k$ satisfont les relations de Hamilton, ce
  qui se ramène à un calcul sur quatre matrices. Le déterminant de l'image vaut
  $|q|^2$ : les quaternions de norme $1$ correspondent donc exactement aux
  matrices de $\mathrm{SU}(2)$.
  \hfill$\blacksquare$
\end{mdframed}

\decsep{}

%=============================================================
%  § 3 — ROTATIONS ET GÉOMÉTRIE
%=============================================================
\secnum{navyblue}{3}{Quaternions et rotations de $\mathbb{R}^3$}

\begin{tcolorbox}[thmbox]
  \textbf{Théorème (Hamilton–Rodrigues).}\enspace%
  \index{quaternions!rotations}\index{Rodrigues (formule de)}
  Soit $\mathbf{n}$ un vecteur unitaire de $\mathbb{R}^3$ et $\theta \in \mathbb{R}$.
  Le quaternion unitaire $q = \cos(\theta/2) + \sin(\theta/2)(n_1 i + n_2 j + n_3 k)$
  représente la rotation d'axe $\mathbf{n}$ et d'angle $\theta$ via :
  \[
    \mathbf{v} \;\longmapsto\; q\,\mathbf{v}\,\bar{q}
  \]
  où $\mathbf{v}$ est identifié au quaternion pur $v_1 i + v_2 j + v_3 k$.
\end{tcolorbox}

\begin{significationintuitive}
  Appliquée à un quaternion pur, la conjugaison $\mathbf{v}\mapsto q\mathbf{v}\bar{q}$
  redonne un quaternion pur et préserve la norme : c'est donc bien une isométrie
  de $\mathbb{R}^3$, et l'on vérifie qu'elle fixe l'axe $\mathbf{n}$. L'apparition
  de l'angle moitié $\theta/2$ vient de ce que $q$ agit \emph{deux fois}, une fois
  à gauche par $q$ et une fois à droite par $\bar{q}$. C'est aussi pourquoi $q$ et
  $-q$ produisent la même rotation, puisque $(-q)\mathbf{v}\,\overline{(-q)} = q\mathbf{v}\bar{q}$ :
  on retrouve ainsi le noyau $\{\pm 1\}$ du recouvrement ci-dessous.
\end{significationintuitive}

\rubriq{Double recouvrement de $\mathrm{SO}(3)$}
L'application $q \mapsto (v \mapsto qv\bar{q})$ est un homomorphisme de groupes
surjectif $\mathrm{Sp}(1) \to \mathrm{SO}(3)$ (où $\mathrm{Sp}(1)$ désigne les
quaternions unitaires), de noyau $\{\pm 1\}$.
Ainsi $\mathrm{SO}(3) \cong \mathrm{Sp}(1)/\{\pm 1\} \cong \mathrm{SU}(2)/\{\pm I\}$.

\decsep{}

%=============================================================
%  § 4 — THÉORÈME DE FROBENIUS ET EXERCICES
%=============================================================
\secnum{exgreen}{4}{Classification, pièges et exercices}

\rubriq{Théorème de classification}

\begin{mdframed}[style=exbar]
  \textbf{Théorème de Frobenius (1877).}\enspace%
  \index{Frobenius (théorème de)}
  Les seules algèbres \emph{associatives} à division de dimension finie sur
  $\mathbb{R}$ sont : $\mathbb{R}$, $\mathbb{C}$ et $\mathbb{H}$.

  \smallskip\noindent
  \textbf{Octonions $\mathbb{O}$.}\enspace%
  \index{octonions}
  En abandonnant aussi l'associativité, on obtient les \emph{octonions} $\mathbb{O}$
  (dimension $8$). Par le théorème de Hurwitz (1898), les seules algèbres
  normées sur $\mathbb{R}$ \emph{dont la norme est multiplicative}
  ($\lVert xy\rVert = \lVert x\rVert\,\lVert y\rVert$) sont
  $\mathbb{R}, \mathbb{C}, \mathbb{H}, \mathbb{O}$
  (dimensions $1, 2, 4, 8$).
\end{mdframed}

\vspace{6pt}
\rubriq{Pièges}

\begin{mdframed}[style=cexbar]
  \textbf{$\mathbb{H}$ n'est pas commutatif.}\enspace%
  Ne pas développer comme dans un anneau commutatif :
  $(p+q)^2 = p^2 + pq + qp + q^2$, et en général $pq + qp \neq 2pq$.
  De même $(pq)^2 \neq p^2q^2$ en général : $(ij)^2 = k^2 = -1$
  alors que $i^2j^2 = (-1)(-1) = 1$.

  \smallskip\noindent
  \textbf{Pas de polynômes bien définis.}\enspace%
  Dans $\mathbb{H}[X]$, l'équation $X^2 + 1 = 0$ a une infinité de solutions :
  tous les quaternions purs unitaires $bi+cj+dk$ avec $b^2+c^2+d^2=1$.
  L'algèbre des polynômes est délicate en non-commutatif.
\end{mdframed}

\vspace{6pt}
\exolabel


\begin{mdframed}[style=exobar]
  \begin{enumerate}
    \item Vérifier les relations $ij=k$, $jk=i$, $ki=j$ à partir de $i^2=j^2=k^2=ijk=-1$.

    \item Montrer que $|pq|=|p||q|$ pour tous $p,q\in\mathbb{H}$.

    \item Montrer que les quaternions unitaires $\mathrm{Sp}(1) = \{q \in \mathbb{H} \mid |q|=1\}$
          forment un groupe (homéomorphe à $\mathbb{S}^3$).

    \item Calculer la rotation correspondant au quaternion $q = \tfrac{1}{2}(1+i+j+k)$.
          Quel est son axe et son angle ?

    \item Montrer que le centre de $\mathbb{H}$ est $\mathbb{R}$, i.e.\ que
          seuls les scalaires commutent avec tous les quaternions.
  \end{enumerate}
\end{mdframed}

\leconfooter{W.\ R.\ Hamilton (16 octobre 1843) — G.\ Frobenius (1877) — A.\ Hurwitz (1898)}
